%% =================================================================
%% Beber et al. IEEE Robotics and Automation Letters, Feb 2026.
%% Reconstruction of page 4 (printed p. 4).
%% Contents: Section III-C continuation (ergodic control defn, Eqs.
%% 4-5, HEDAC practical implementation); Section III-D Trajectory
%% Execution (III-C-1 Online Elasticity/Viscosity Estimation including
%% Eq. 6 the DRM/Hertzian viscoelastic contact model, III-C-2 Filter
%% Design the EKF); Sections III-E Target Distribution Update and
%% III-F Search Termination Criterion; start of Section IV Simulation
%% Results (IV-A Setup and Baselines).
%% =================================================================

\documentclass[11pt]{article}

\usepackage[margin=0.85in]{geometry}
\usepackage{amsmath, amssymb}
\usepackage{mathrsfs}
\usepackage{xcolor}
\usepackage{fancyhdr}

\pagestyle{fancy}
\fancyhf{}
\lhead{\small 4\ \ IEEE RAL, Accepted Feb 2026}
\rhead{\small }
\cfoot{\small 4 $\to$ \arabic{page}}
\renewcommand{\headrulewidth}{0.4pt}

\setcounter{page}{1}
\setcounter{equation}{3}  %% Eqs. 1-3 landed on p. 3; p. 4 starts with (4)
\setlength{\parskip}{6pt}
\setlength{\parindent}{0pt}

\begin{document}

% ---------- Continuation of Section III-C ergodic control defn ----------
where $\phi_{\sigma}$ is a Gaussian kernel. The coverage error with
respect to a target distribution $m(x)$ is defined as the spatial
diffusion of $m(x)$ through the RBF kernel compared with the coverage
density
\begin{equation}
e(x, t) = (\phi_{\sigma} * m)(x) - c(x, t),
\end{equation}
To avoid local minima and improve spatial coordination, HEDAC
smooths this error using a stationary heat equation:
\[
\varrho\, \Delta u(x, t) = \beta\, u(x, t) + \gamma\, a(x, t) - s(x, t),
\qquad
\dot{z} = v_{a} \cdot \frac{\nabla u}{\|\nabla u\|},
\]
where $s$ highlights under-sampled regions, and $a$ ensures
collision avoidance between the agents.

\textbf{2) Practical Implementation:} Different from [19], this work
implemented the non-stationary diffusion equation, as was done in
[26], for a more locally consistent exploration behaviour. This means
that $\dot{u} \ne 0$ is defined as:
\begin{equation}
\dot{u}(x, \tau) = \Delta u(x, \tau),
\end{equation}
allowing us to control the desired smoothness. In this formulation,
an initial condition is required: $u(x, 0) = s(x, 0)$. During the
experiments, a second-order agent is used to generate the trajectory.
The trajectory is updated with a frequency $\Delta T$ chosen to
ensure smooth motion and adequate control loop resolution (typically
100 Hz in our implementation, 5 times slower than the low-level
controller). The planning phase respects constraints on maximum
velocity to ensure physical plausibility during data collection
imposed by the second-order dynamics inside the ergodic controller.

% ---------- Section III-D ----------
\subsection*{D.\ Trajectory Execution and Data Collection}

While the trajectory is executed, tissue elasticity is estimated
continuously using an Extended Kalman Filter (EKF). Data are sampled
at a rate that depends on the robot motion and the indenter geometry.
Specifically, higher velocities or smaller indenters require denser
sampling to avoid aliasing effects in the stiffness map.

\textbf{1) Online Elasticity and Viscosity Estimation:} Accurately
estimating the mechanical response of soft tissues in real time is a
challenging task, as it requires modelling the contact interaction
between the probe and the tissue surface and inferring tissue
parameters from force measurements and robot motion. In this work,
we adopt a viscoelastic contact model to improve the fidelity of
force prediction during dynamic palpation, while subsequent
abnormality detection focuses on the estimated elastic stiffness.
To achieve this, we adopt DRM [27], which allows mapping the
three-dimensional contact mechanics into an equivalent
two-dimensional representation. \pdfcomment[color=green,subject={},open=false,opacity=0.7]{User input: that the quantities inferred from the measurements use a mathematical model that ONLY ACCOUNTS FOR THE LOCAL PHYSICS in the ergodic paper\\Claude comments: EVIDENTIAL RELATIONSHIP: supports. Hertzian contact mechanics + viscoelastic damping = purely local physics. Eq.~6 takes only local point-contact variables ($d, \dot{d}, E_f, \eta, \nu, R$) --- no bulk PDE, no global FEM solve, no neighbour coupling. The Hertzian half-space assumption is the definitive local-physics fingerprint. Contrast with MBT which solves a global elastic PDE.}\pdfcomment[color=green,subject={},open=false,opacity=0.7]{User input: that the inference model treats forces as NORMAL, i.e. no TANGENTIAL component\\Claude comments: EVIDENTIAL RELATIONSHIP: supports. $F_{TOT}$ is a scalar (not vector) --- no normal/tangential decomposition. All inputs ($d, \dot{d}, R, E_f, \eta, \nu$) are rotationally symmetric about the contact normal. No friction coefficient $\mu$ or shear modulus $G$ appears. The EKF state has no tangential kinematic variables. Hertzian contact mechanics is the canonical frictionless-normal-load theory.}\sethlcolor{hlgreen}\hl{Under this framework, the force
exerted by the tissue during indentation is a nonlinear function of
penetration depth and velocity:}
\begin{equation}
F_{\mathrm{TOT}}(d, \dot{d}) =
\frac{4}{3}\, \frac{E_{f}}{1 - \nu^{2}}\, \sqrt{R\, d}\, d
+ \frac{4}{1 - \nu}\, \eta\, \sqrt{R\, d}\, \dot{d},
\end{equation}
\sethlcolor{hlgreen}\hl{where $d$ is the penetration depth, $\dot{d}$ its derivative
(velocity), $E_{f}$ is the Young's modulus of the material, $\eta$
is the viscosity coefficient, $\nu$ is the Poisson's ratio, and $R$
is the radius of the spherical indenter [18]. The first term models
elastic deformation according to Hertzian contact mechanics, while
the second term accounts for viscoelastic damping.}

\textbf{2) Filter Design:} \pdfcomment[color=cyan,subject={},open=false,opacity=0.7]{User input: what was inferred from a given SINGLE force measurement?\\Claude comments: A single force sample drives one Bayesian update of an 8-dim EKF state: penetration depth $d$, velocity $\dot{d}$, local stiffness $k$, local damping $\lambda$, and their time derivatives. The contact-mechanics model (Eq.~6) relates $F$ to $(d, \dot{d}, E_f, \eta, \nu, R)$ --- the EKF inverts it online. $d$ itself is not directly measured; only $F$ is.}\sethlcolor{hlcyan}\hl{The main challenge lies in the fact that
$d$ is not directly measurable during robotic interaction. Therefore,
we use an EKF to estimate the system state, which includes both $d$
and $\dot{d}$, as well as the unknown tissue parameters. The system
is modelled using a discrete-time nonlinear state-space
representation, as detailed in [18]. The state vector is defined as}
\[
[\,d,\ k,\ \lambda,\ \dot{d},\ \dot{k},\ \dot{\lambda},\ \ddot{k},\
\ddot{\lambda}\,]^{\top},
\]
\sethlcolor{hlcyan}\hl{where $\kappa$ and $\lambda$ correspond to the local stiffness and
damping parameters, respectively. The filter assumes as input the
measured contact force and comprises the effective interaction mass.
Process and measurement noise are also included to account for model
uncertainties and sensor imperfections. This filtering approach
allows real-time estimation of the viscoelastic parameters even
during continuous contacts with indentation speeds up to 5 mm/s
[28], being suitable for dynamic exploration tasks.}

% ---------- Section III-E ----------
\subsection*{E.\ Target Distribution Update}

Once a sufficient number of new samples are acquired, the GPR model
is updated to refine the target stiffness distribution. Since GPR
inference has cubic time complexity with respect to the number of
samples, we avoid updating the model at every time step. Instead,
updates are performed at a fixed frequency (e.g., 1 Hz), allowing
batch accumulation of new data while keeping computation time
bounded.

% ---------- Section III-F ----------
\subsection*{F.\ Search Termination Criterion}

Different from existing approaches that rely on predefined stopping
conditions [15], or fixed time/step limits [16], we employ an
information-theoretic stopping criterion on $\alpha$, which is based
on the ergodic metric. The value of $\alpha$ governs the trade-off
between exploration and exploitation. Larger values of $\alpha$
correspond to an early exploration phase and favour earlier
termination, but may lead to incomplete spatial coverage. Conversely,
smaller values of $\alpha$ indicate that the exploration has largely
converged toward the target EID distribution, resulting in extended
exploitation phases, longer execution times, and only marginal
refinement of already identified regions. In practice, the stopping
threshold on $\alpha$ is selected within the exploitation-dominated
regime, such that further reductions correspond to diminishing
returns in reconstruction accuracy. This choice is supported by
empirical observations indicating saturation of the RMSE as $\alpha$
approaches a plateau, and reflects a practical balance between
coverage completeness, estimation accuracy, and exploration duration.
The advantage of using this metric lies in its ability to quantify
how thoroughly the area has been explored, without relying on prior
knowledge such as fixed time limits or the number and characteristics
of stiff regions that need to be found. This criterion is preferable
to a fixed-time limit, as it adapts the exploration duration to the
complexity of the stiffness distribution and the desired level of map
details, avoiding both premature termination and unnecessary probing.

% ---------- Section IV ----------
\section*{IV.\ Simulation Results}

\subsection*{A.\ Setup and Baselines}

The proposed search algorithm was evaluated on synthetic stiffness
distributions representative of biological soft tissues.

\end{document}
